\documentclass[border=10pt]{standalone}

% пакеты для математики
\usepackage{amsmath,amsfonts,amssymb,amsthm,mathtools}

\usepackage[british,russian]{babel} % выбор языка для документа

% Математические шрифты
\usepackage{unicode-math}
\setmathfont[math-style=upright]{[Neo Euler.otf]}

\usepackage{tikz}
\usepackage{xcolor}

\usetikzlibrary{calc,shapes,positioning}

\begin{document}
    \begin{tikzpicture}[line cap=round,line join=round,x=2.0cm,y=2.0cm]
        \draw [->, line width=1.pt] (-1.3,0) -- (3.5,0) node[below] {\small $x$};
        \draw [->, line width=1.pt] (0,-1.4) -- (0,1.5) node[left] {\small $g(x)$};

        % пунктир к вершине
        \draw [line width=1.pt,dash pattern=on 3pt off 3pt] (0,1) -- (1,1) -- (1,0);
        \draw (0,1) node[left] {\small $1$};
        \draw (1,0) node[below] {\small $1$};
        \draw (2,0) node[above] {\small $2$};
        \draw (0,0) node[below left] {\small $0$};

        % ломаная: хвост уезжает вниз
        \draw [line width=2.pt] (-1.2,0) -- (0,0) -- (1,1) -- (3.2,-1.2);

        % изломы
        \foreach \p in {(0,0), (1,1)}{
            \draw [line width=1.pt,fill=white] \p circle (2.5pt);
        }
    \end{tikzpicture}
\end{document}
